\chapter{Trazabilidad Histórica de la Revisión}
\label{ap:original_adaptado}

Este apéndice documenta cómo evolucionó el documento hasta la versión final. No incluye tablas antiguas como evidencia principal, porque algunas pertenecen a fases exploratorias o a marcos de análisis que fueron reemplazados durante la revisión. Su función es dejar constancia de la evolución metodológica y de las decisiones de cierre.

\section{De resultados aislados a trayectoria experimental}

La versión inicial del trabajo enfatizaba rankings y comparaciones agregadas. Durante la revisión se reorganizó el Capítulo~\ref{ch:resultados} como \textit{Experimentos y Resultados}, incorporando selección preliminar, reducción de candidatos, optimización intermedia, comparación final, señales representativas, PSD, análisis de contribución temporal, costo computacional y frontera práctica. Este cambio evita que el lector vea la comparación TRSS--MNE como una elección directa sin justificación previa.

\section{Cambio de marco interpretativo}

El cierre final se estructura alrededor de preguntas de investigación, comparaciones pareadas descriptivas, tasas de victoria, intervalos de incertidumbre y criterios prácticos de decisión. Los contrastes inferenciales no son el eje del documento final. Esta decisión es coherente con el objetivo de la tesis: caracterizar cuándo TRSS es útil, cuándo MNE sigue siendo preferible y qué condiciones limitan la generalización.

\section{Correcciones de lenguaje y formato}

La revisión final eliminó términos poco adecuados para el registro académico del documento, aclaró la expresión “métodos espaciales de referencia”, retiró menciones a archivos locales, redujo guiones largos innecesarios y ajustó tablas cuya información era más clara en prosa. También se rediseñaron figuras y tablas cuando el formato visual interfería con la lectura.

\section{Métodos invalidados o no concluyentes}

Algunas variantes experimentales quedaron fuera del cierre final. Visibility NNK se excluyó por error de implementación. Grafos aprendidos, grafos adaptativos, correlación funcional y métodos temporales alternativos permanecen como evidencia exploratoria o trabajo futuro cuando no cumplen condiciones de costo, reproducibilidad o independencia metodológica suficientes para sostener conclusiones finales.

\section{Uso de este apéndice}

Este apéndice no debe leerse como una fuente adicional de resultados cuantitativos. Para resultados finales deben consultarse el Capítulo~\ref{ch:resultados}, el Capítulo~\ref{ch:conclusiones} y la síntesis extendida del Apéndice~\ref{ap:tablas_extensas}. Su propósito es documentar la trazabilidad editorial y metodológica de la versión final.
